\subsection{希尔伯特–塞缪尔函数的次数}

定理 3. --- 设 A 为诺特局部环，q 为不同于 A 的 A 的理想，M 为有限生成 A-模，且 M/qM 具有有限长度。则整数 \(d_{\mathfrak{q}}(\mathbf{M})\) 是 A-模 M 的维数（§ 1, no. 4, def. 8）。

不妨设 \(\mathbf{M} \neq 0\)。先证明不等式 \(d_{\mathfrak{q}}(\mathbf{M}) \leqslant \dim_{\mathbf{A}}(\mathbf{M})\)。由 § 3, no. 5 的 prop. 9 的 cor. 2，存在 \(x_1, ..., x_r \in \mathfrak{q}\)，其中 \(r = \dim_{\mathbf{A}}(\mathbf{M})\)，且 \(\text{long}(\mathbf{M}/\sum_{i=1}^{r} x_i \mathbf{M}) < +\infty\)；置 \(\mathfrak{x} = \sum_{i=1}^{r} x_i \mathbf{A}\)。由 no. 3 的 prop. 4，有 \(d_{\mathfrak{x}}(\mathbf{M}) \leqslant r\)；又有 \(\mathfrak{x} \subset \mathfrak{q}\)，因而 \(\mathrm{H}_{\mathbf{M},\mathfrak{q}}^{(1)} \leqslant \mathrm{H}_{\mathbf{M},\mathfrak{x}}^{(1)}\)，于是（no. 1 的引理 2）

\[
d _ {\mathfrak{q}} (\mathbf {M}) \leqslant d _ {\mathfrak{x}} (\mathbf {M}) \leqslant r = \dim_ {\mathrm{A}} (\mathbf {M}).
\]

现在对 \(\dim_{\mathbf{A}}(\mathbf{M})\) 归纳证明不等式 \(\dim_{\mathbf{A}}(\mathbf{M}) \leqslant d_{\mathfrak{q}}(\mathbf{M})\)；当 \(\dim_{\mathbf{A}}(\mathbf{M}) = 0\) 时显然成立。

假设 \(\dim_{\mathbf{A}}(\mathbf{M}) > 0\)，并且 \(\dim_{\mathbf{A}}(\mathbf{N}) \leqslant d_{\mathfrak{q}}(\mathbf{N})\) 对每个满足 \(\dim_{\mathbf{A}}(\mathbf{N}) < \dim_{\mathbf{A}}(\mathbf{M})\) 的有限生成 A-模 N 都成立。若 \(0 = \mathbf{M}_0 \subset \mathbf{M}_1 \subset \ldots \subset \mathbf{M}_n = \mathbf{M}\) 是 M 的合成序列，则有 \(\dim_{\mathbf{A}}(\mathbf{M}) = \sup(\dim_{\mathbf{A}}(\mathbf{M}_i / \mathbf{M}_{i-1}))\)（§ 1, no. 4, prop. 9），且 \(d_{\mathfrak{q}}(\mathbf{M}) = \sup(d_{\mathfrak{q}}(\mathbf{M}_i / \mathbf{M}_{i-1}))\)（No. 3, prop. 5）。根据 IV, § 1, no. 4 的 th. 1，因此可假设 M 形如 A/p，其中 p 是 A 的素理想；有 \(\mathfrak{p} \ne \mathfrak{m}_{A}\)，因为 \(\dim_{\mathbf{A}}(\mathbf{M}) > 0\)。取 \(x \in \mathfrak{m}_{\mathbf{A}} - \mathfrak{p}\)；M = A/p 上的倍乘映射 \(x_{\mathbf{M}}\) 是单射，并且有正合序列

\[
0 \longrightarrow \mathrm{M} \xrightarrow {x _ {\mathrm{M}}} \mathrm{M} \longrightarrow \mathrm{M} / x \mathrm{M} \longrightarrow 0.
\]

由 § 3, no. 2 的 prop. 3，有 \(\dim_{\mathbf{A}}(\mathbf{M}/x\mathbf{M}) = \dim_{\mathbf{A}}(\mathbf{M}) - 1\)；由 no. 3 的 prop. 5 以及前面的正合序列，有 \(d_{\mathfrak{q}}(\mathbf{M}/x\mathbf{M}) \leqslant d_{\mathfrak{q}}(\mathbf{M}) - 1\)。根据归纳假设，有

\[
\dim_ {\mathbf {A}} (\mathbf {M}) = \dim_ {\mathbf {A}} (\mathbf {M} / x \mathbf {M}) + 1 \leqslant d _ {\mathfrak {q}} (\mathbf {M} / x \mathbf {M}) + 1 \leqslant d _ {\mathfrak {q}} (\mathbf {M}),
\]

这就完成了证明。

推论。--- 设 A 为诺特环，M 为有限生成 A-模，q 是 A 的理想，且 M/qM 具有有限长度。则 \(d_{\mathfrak{q}}(\mathbf{M})\) 是各维数 \(\dim_{\mathrm{A}_{\mathfrak{m}}}\left(\mathrm{M}_{\mathfrak{m}}\right)\) 的上确界，其中 \(\mathfrak{m}\) 遍历有限集 \(S = \operatorname{Supp}(\mathbf{M}) \cap \mathrm{V}(\mathfrak{q})\)；而 \(e_{\mathfrak{q}}(\mathbf{M})\) 是将 \(e_{\mathfrak{q}_{\mathfrak{m}}}\left(\mathrm{M}_{\mathfrak{m}}\right)\) 对满足 \(\mathfrak{m}\in S\) 且 \(\dim_{\mathrm{A}_{\mathfrak{m}}}\left(\mathrm{M}_{\mathfrak{m}}\right) = d_{\mathfrak{q}}(\mathbf{M})\) 的项相加所得的和。

对于每个整数 \(n\)，\(\mathbf{M} / \mathfrak{q}^n\mathbf{M}\) 的长度等于 \(\mathrm{long}_{\mathbf{A}_{\mathfrak{m}}}(\mathbf{M}_{\mathfrak{m}} / \mathfrak{q}_{\mathfrak{m}}^n\mathbf{M}_{\mathfrak{m}})\) 之和（IV, § 2, no. 5, cor. 1 to prop. 7 and corollary to prop. 8）。因此，有 \(\mathbf{H}_{\mathbf{M},\mathfrak{q}} = \sum_{\mathfrak{m}\in S}\mathbf{H}_{\mathbf{M}_{\mathfrak{m}},\mathfrak{q}_{\mathfrak{m}}}\)，从而得到推论。

备注。--- 1) 我们还有 \(d_{\mathfrak{q}}(\mathbf{M}) = \sup_{\mathfrak{m} \in \mathbf{V}(\mathfrak{q})} \dim(\mathbf{M}_{\mathfrak{m}})\)，也就是说，\(d_{\mathfrak{q}}(\mathbf{M}) = \dim(\hat{\mathbf{M}})\)，其中 \(\hat{\mathbf{M}}\) 是 \(\mathbf{M}\) 关于 q-进拓扑的完备化（§ 3, No. 4, prop. 8）。

\begin{enumerate}
\def\labelenumi{\arabic{enumi})}
\setcounter{enumi}{1}
\tightlist
\item
  假设 q 包含于 A 的根基中；则 \(\dim(\tilde{\mathbf{M}})=\dim(\mathbf{M})\)（loc. cit., cor. 1），因而 \(d_{\mathfrak{q}}(\mathbf{M})=\dim(\mathbf{M})\)。
\end{enumerate}

